\section*{Chapter \ref{ch:b2:bayesian-games} --- Games of Incomplete Information}

\begin{solution}{exo:b2:bayesian-games:1}
% code: exercises/solutions.py first_price_bid
In the second-price auction the bidder with value 80 wins and pays 50. In the first-price equilibrium each bids two thirds of
her value: 20, 33.3 and 53.3; the bidder with value 80 wins and pays 53.3.
\end{solution}

\begin{solution}{exo:b2:bayesian-games:2}
% figdata: bids_summary
All eighteen with three bidders. With six, fourteen bid above, one at and three below the prediction. The average bids
were 0.81 and 0.87 per dollar of value, against predictions of 0.67 and 0.83.
\end{solution}

\begin{solution}{exo:b2:bayesian-games:3}
% code: exercises/solutions.py revenue
$(n - 1)/(n + 1)$: $1/3 = 0.333$ with 2 bidders and $9/11 = 0.818$ with 10. More bidders raise the expected highest value and
reduce shading.
\end{solution}

\begin{solution}{exo:b2:bayesian-games:4}
Bids systematically above values, which lose money whenever the second-highest bid falls between the value and the
bid, and which persist with experience, refute it. Bids that differ from values only by amounts that rarely change who
wins or what she pays leave it standing: bidding one's value is only weakly dominant, and small deviations cost little.
\end{solution}

\begin{solution}{exo:b2:bayesian-games:5}
% code: exercises/solutions.py cournot_private
Firm 1's \omterm{def:b2:static-games:best-response}{best response} to the expected output of firm 2 with the solutions for $q_2^H$ and $q_2^L$ gives $q_1 = (12 - 4 + 0.25 \times 4 +
0.75 \times 2)/3 = 3.5$. Then $q_2^H = (8 - 3.5)/2 = 2.25$ and $q_2^L = (10 - 3.5)/2 = 3.25$.
\end{solution}

\begin{solution}{exo:b2:bayesian-games:6}
A bidder with value 10 faces one rival who bids 4. Bidding 10 wins and earns $10 - 10 = 0$; bidding 5 also wins and earns 5. So
bidding one's value is not a \omterm{def:b2:static-games:best-response}{best response} to every strategy of the others.
\end{solution}

\begin{solution}{exo:b2:bayesian-games:7}
% code: exercises/solutions.py bayes, bayes_poor_report
Reported rich: $0.3 \times 0.8/(0.3 \times 0.8 + 0.7 \times 0.2) = 0.24/0.38 = 0.632$. Reported poor: $0.3 \times 0.2/(0.3 \times 0.2 + 0.7 \times 0.8) = 0.06/0.62 =
0.097$.
\end{solution}

\begin{solution}{exo:b2:bayesian-games:8}
% code: exercises/solutions.py jar
% ledger: b2:thaler1988-jar
$10.01 - 8 = 2.01$ dollars. The winner is not the average bidder but the one with the highest estimate; with many bidders, the
highest of the estimates lies far above their average, enough to exceed the true value although the average fell short of
it.
\end{solution}

\begin{solution}{exo:b2:bayesian-games:9}
% code: exercises/solutions.py first_price_bid, risk_aversion_from_slope
$2/(3 - 0.5) = 0.8$. A slope of 0.75 requires $3 - r = 2/0.75 = 2.667$, so $r = 0.333$.
\end{solution}

\begin{solution}{exo:b2:bayesian-games:10}
Suppose the others bid $kv$. A bid $b \le k\bar v$ beats one rival with probability $b/(k\bar v)$, so the bidder maximises $(v - b)\big(b/(k\bar
v)\big)^{n-1}$. The first-order condition $(n - 1)(v - b) = b$ gives $b = (n - 1)v/n$, so $k = (n - 1)/n$. The objective is zero at $b = 0$ and $b = v$ and
rises then falls between them, so the condition picks the maximum; a bid above $k\bar v$ wins for sure, as does $k\bar v$, which costs
less; so no \omterm{def:b2:bayesian-games:type}{type} gains from any deviation.
\end{solution}

\begin{solution}{exo:b2:bayesian-games:11}
The highest of $n$ values is below $x$ with probability $x^n$, so its density is $nx^{n-1}$ and its expectation $\int_0^1 nx^n\,dx = n/(n +
1)$. The second-highest is below $x$ when at most one value exceeds $x$: probability $x^n + nx^{n-1}(1 - x)$, with density $n(n - 1)x^{n-2}(1 -
x)$ and expectation $n(n - 1)\big(1/n - 1/(n + 1)\big) = (n - 1)/(n + 1)$. The first-price auction yields $(n - 1)/n$ times the highest value, with
expectation $(n - 1)/n \times n/(n + 1) = (n - 1)/(n + 1)$; the second-price auction yields the second-highest value.
\end{solution}

\begin{solution}{exo:b2:bayesian-games:12}
Write $\varepsilon_i = \bar\varepsilon(2u_i - 1)$ with $u_i$ uniform on $[0, 1]$. The highest error is $\bar\varepsilon(2u_{\max} - 1)$, with expectation $\bar\varepsilon(2n/(n + 1) - 1) =
\bar\varepsilon(n - 1)/(n + 1)$.
\end{solution}

\begin{solution}{pb:b2:bayesian-games:1}
% code: exercises/solutions.py licence
\textbf{Part I.}
1. Its value for the licence.
2. $1/4$.
3. $800 \times 4/5 = 640$.
4. $800 \times 3/5 = 480$.
5. $(600/800)^3 = 0.422$.
6. Each firm's value depends on its own boats, costs and markets, and would not change if it learned the others' values.
\textbf{Part II.}
7. $b(v) = 3v/4$.
8. 450.
9. $(600 - 450) \times 0.422 = 63.28$.
10. $3/4 \times 640 = 480$.
11. $3/(4 - 0.5) \times 600 = 514.3$.
12. Bidding its value would earn nothing even when it wins; shading lowers the chance of winning but raises the gain.
\textbf{Part III.}
13. Its value.
14. The expected second-highest value, 480.
15. If it wins, it pays the highest of three values below 600, on average $600 \times 3/4 = 450$; its expected payoff is $(600 - 450) \times
0.422 = 63.28$, as in Part II.
16. The first-price auction: same expected revenue, less risk.
17. The first-price auction: \omterm{def:b2:expected-utility:risk-aversion}{risk aversion} raises first-price bids and leaves second-price bids at values.
18. The second-price auction: bidding one's value is best whatever the others do, so a firm need not guess their values
or strategies, and mistakes are less likely.
\textbf{Part IV.}
19. 500.
20. $500 - 100 \times 3/5 = 440$.
21. Winning means its signal was the highest, so the licence is worth 440 in expectation, 60 less than the 500 it assumed:
it can pay up to 60 more than the licence is worth.
22. $100 \times 7/9 = 77.8$.
23. With more firms, the highest of the signals is more likely to carry a large positive error.
24. Both formats raise an expected 480 thousand dalens; with a common value and four firms, a firm should lower its estimate
by 60 thousand dalens, the expected error of the highest signal.
\end{solution}
